\documentclass[10pt,a4paper]{amsart}
\usepackage[margin=19mm]{geometry}
\usepackage{amsmath,amssymb,mathtools}
\usepackage[scr=rsfs,cal=euler]{mathalfa}
\usepackage{xfrac}
\usepackage[unicode,hidelinks,bookmarks=false]{hyperref}
\newcommand{\ie}{i.e.\ }
\newcommand{\cf}{cf.\ }
\newcommand{\resp}{respectively\ }
\newcommand{\Subsection}{\subsection}
\providecommand{\nobreakdash}{\nobreak\hbox{-}}
\setlength{\emergencystretch}{2em}
\multlinegap0pt
\usepackage{amssymb}
\usepackage[shortlabels]{enumitem}
\newtheorem{theorem}{Theorem}
\newcommand{\Lap}[1]{\Delta_{#1}}
\newcommand{\N}{\mathbb{N}}
\newcommand{\supp}{\operatorname{supp}}
\title{An improved strong comparison principle for singular $p$-Laplace equations, with applications}
\author{Phuong Le}
\date{Source: arXiv:2609.07877v1 (CC BY 4.0)}
\begin{document}
\maketitle

\noindent\emph{Source and licence.} An improved strong comparison principle for singular $p$-Laplace equations, with applications. Version arXiv:2609.07877v1 (CC BY 4.0), distributed by arXiv under the Creative Commons Attribution 4.0 International licence, http://creativecommons.org/licenses/by/4.0/, as recorded in the arXiv licence metadata for this exact version and shown on the abstract page https://arxiv.org/abs/2609.07877v1. Full licence notice: \texttt{licenses/arxiv-2609.07877-CC-BY-4.0.txt}.\par\smallskip

\noindent\textit{Editorial scope.} Counted unit: the article's theorem thm:abstract (source0.tex lines 326--335). The article's complete original proof of it is delivered with it (source0.tex lines 942--1217, one \texttt{proof} environment of 276 lines, which the article places away from the statement). No mathematics is written, repaired or re-derived: the statement and the proof are the article's own lines, in the article's own notation, and no retained line is substituted or re-flowed.  Retained uncounted as the source-local context the proof needs: lines 237--239; lines 807--824; lines 819--822; lines 895--897; lines 899--901.  Nothing was omitted from any retained span, so the package carries no omission record.  No retained line contains a citation, so the wrapper carries no bibliography and no bibliography entry.  No retained line contains a figure environment, an image-inclusion command or drawing code, so no image is delivered and the package declares no asset dependency.  Editorial formatting only, in addition: an AMS \texttt{amsart} preamble that restates the article's own declarations and macros used by the retained lines, namely the environments \texttt{theorem} on separate counters, together with the standard packages those lines need.  The retained environments print in this document as Theorem 1 in order of appearance, whereas the article prints the same items with the numbering of its own full document.  The complete native source of the article as distributed in the arXiv e-print for this exact version is bundled unchanged as \texttt{sources/209668/source0.tex}.
\begin{equation}\label{eq:scp-hyp}
-\Lap{p}(u)+\Lambda u\ \le\ -\Lap{p}(v)+\Lambda v,\qquad u\le v\quad\text{in }\Omega
\end{equation}

\begin{theorem}\label{thm:sobolev}
Assume $\rho$ satisfies $(\mathcal M_\kappa)$ in $\Omega'$ and set
\begin{equation}\label{eq:qdef}
q:=\begin{cases}\ \dfrac{2(N-2+\kappa)}{N-2}=2+\dfrac{2\kappa}{N-2}, & N\ge3,\\[2mm]
\ 2+\kappa, & N=2 .\end{cases}
\end{equation}
Then $q>2$ and there exists $C_S=C_S(N,\kappa,\lambda_0,\Lambda_0,\Omega')>0$ such that
\begin{equation}\label{eq:sobolev}
\|w\|_{L^{q}(\Omega',\mu)}\ \le\ C_S\,\|Dw\|_{L^2(\Omega',\mu)}
\qquad\forall\,w\in H^{1,2}_{0,\rho}(\Omega') .
\end{equation}
Moreover, for every ball $B\subseteq\Omega'$ and every $w\in H^{1,2}_{\rho}(B)$,
\begin{equation}\label{eq:sobolev-mean}
\|w-\overline w_B\|_{L^{q}(B,\mu)}\ \le\ C_S'\,\|Dw\|_{L^2(B,\mu)},\qquad
\overline w_B:=\frac1{|B|}\int_Bw\,dx ,
\end{equation}
with $C_S'$ depending in addition on $\operatorname{diam}(\Omega')$ but not on $B$.
\end{theorem}

\begin{equation}
\|w-\overline w_B\|_{L^{q}(B,\mu)}\ \le\ C_S'\,\|Dw\|_{L^2(B,\mu)},\qquad
\overline w_B:=\frac1{|B|}\int_Bw\,dx ,
\end{equation}

\begin{equation}\label{eq:cacc-comp}
\int\rho\,\eta^2|D\widetilde w|^2\,dx\ \le\ \dot C\,r^2\int \widetilde w^{\,2}\big(\eta^2+\rho|D\eta|^2\big)\,dx ,
\end{equation}

\begin{equation}\label{eq:cacc-log}
\int\rho\,\eta^2|D\widetilde w|^2\,dx\ \le\ \dot C\int\big(\eta^2+\rho|D\eta|^2\big)\,dx ,
\end{equation}

\begin{theorem}[Abstract weak Harnack comparison inequality]\label{thm:abstract}
Let $u,v\in C^1(\Omega)$ with $u\le v$ and assume \eqref{eq:scp-hyp}. Let $\Omega'\Subset\Omega$ and
suppose that the weight $\rho:=(|Du|+|Dv|)^{p-2}$ satisfies $(\mathcal M_\kappa)$ in $\Omega'$ for
some $\kappa>0$. Then there exists $\chi>1$ such that, whenever
$\overline{B_{6\delta}(x)}\subset\Omega'$ and $0<s<\chi$, there is a constant $C>0$ with
\[
\Big(\int_{B_{2\delta}(x)}(v-u)^s\,d\mu\Big)^{1/s}\le C\,\inf_{B_{\delta}(x)}(v-u).
\]
Consequently the strong comparison principle holds in $\Omega'$.
\end{theorem}

\begin{proof}[Proof of Theorem \ref{thm:abstract}]
Let $\rho=(|Du|+|Dv|)^{p-2}$ satisfy $(\mathcal M_\kappa)$ in $\Omega'$ and let $q>2$ and $C_S$ be
as in Theorem \ref{thm:sobolev}. Set $\chi:=q/2>1$.

\emph{Step 1: the basic iteration inequality in $\mu$.} Since $\rho\ge\lambda_0$ we have, for any
$\eta\in C_c^1$ with $\|D\eta\|_\infty\le K$,
\begin{equation}\label{eq:passtomu}
\int \widetilde w^{2}\big(\eta^2+\rho|D\eta|^2\big)dx
\le \big(\lambda_0^{-1}+1\big)\int \widetilde w^{2}\big(\eta+|D\eta|\big)^2\rho\,dx
=\big(\lambda_0^{-1}+1\big)\big\|\widetilde w(\eta+|D\eta|)\big\|_{L^2(\mu)}^2 .
\end{equation}
On the other hand, $\eta\widetilde w\in H^{1,2}_{0,\rho}$ and, by Theorem \ref{thm:sobolev},
\[
\|\eta\widetilde w\|_{L^{q}(\mu)}^2\le C_S^2\int\rho|D(\eta\widetilde w)|^2dx
\le 2C_S^2\Big[\int\rho\eta^2|D\widetilde w|^2+\int\rho\widetilde w^2|D\eta|^2\Big].
\]
Combining with \eqref{eq:cacc-comp} and \eqref{eq:passtomu} we obtain the fundamental inequality
\begin{equation}\label{eq:fundamental}
\|\eta\widetilde w\|_{L^{q}(\mu)}^2\ \le\ \ddot C\,(1+|r|)^2\,\big\|\widetilde w\,(\eta+|D\eta|)\big\|_{L^2(\mu)}^2 ,
\end{equation}
where the summand $1$ in $(1+|r|)^2$ accounts for the term $\int\rho\widetilde w^2|D\eta|^2$, which
does not carry the factor $r^2$ produced by \eqref{eq:cacc-comp}.
Analogously, from \eqref{eq:cacc-log},
\begin{equation}\label{eq:fundamental-log}
\int\rho|D\widetilde w|^2\eta^2dx\le \ddot C\int(\eta^2+\rho|D\eta|^2)dx\le \ddot C\,\mu(\Omega')\big(\lambda_0^{-1}+K^2\big) .
\end{equation}

\emph{Step 2: the reverse-H\"older chain.} Define, for $s\ne0$ and $R>0$,
\[
\Phi(s,R,z):=\Big(\int_{B_R(x)}|z|^{s}\,d\mu\Big)^{1/s},\qquad
\Phi(+\infty,R,z):=\sup_{B_R(x)}|z|,\qquad \Phi(-\infty,R,z):=\inf_{B_R(x)}|z| .
\]
Since $\lambda_0\,dx\le d\mu$ and $\mu(\Omega')<\infty$, the measure $\mu$ is finite and mutually
absolutely continuous with respect to $dx$; hence the two limits
$\lim_{s\to\pm\infty}\Phi(s,R,z)=\Phi(\pm\infty,R,z)$ hold, and H\"older's inequality reads
\begin{equation}\label{eq:holderPhi}
\Phi(s_1,R,z)\le\mu(B_R(x))^{\frac1{s_1}-\frac1{s_2}}\,\Phi(s_2,R,z),\qquad 0<s_1<s_2 .
\end{equation}

Fix $\delta\le h'<h''\le5\delta$ and let $\eta\in C_c^1(B_{h''}(x))$ with $\eta\equiv1$ on
$B_{h'}(x)$, $0\le\eta\le1$ and $|D\eta|\le\frac{2}{h''-h'}$. Since $h''-h'\le5\delta$, we have
$\eta+|D\eta|\le\frac{C(\delta)}{h''-h'}$. With $\beta<0$, $\beta\ne-1$, $r:=\beta+1<1$ and
$\widetilde w=w_\varsigma^{r/2}$ we have $\widetilde w^{\,2}=w_\varsigma^{r}$, so that
$\|\widetilde w\|_{L^2(B_{h''}(x),\mu)}^2=\Phi(r,h'',w_\varsigma)^{r}$ and
\[
\|\eta\widetilde w\|^2_{L^{q}(\mu)}
\ \ge\ \Big(\int_{B_{h'}(x)}w_\varsigma^{\chi r}\,d\mu\Big)^{2/q}
=\Phi(\chi r,h',w_\varsigma)^{r} ,
\]
because $\frac{q r}{2}=\chi r$ and $\frac2q=\frac1\chi$. Therefore \eqref{eq:fundamental} becomes
\begin{equation}\label{eq:chain}
\Phi(\chi r,h',w_\varsigma)^{r}\ \le\ \Big(\frac{C(1+|r|)}{h''-h'}\Big)^{2}\,\Phi(r,h'',w_\varsigma)^{r},
\end{equation}
with $C=C(p,\Lambda,\delta,q,\lambda_0,\Lambda_0,\|v\|_{L^\infty},\|Du\|_{L^\infty},\|Dv\|_{L^\infty})$.
Raising \eqref{eq:chain} to the power $1/r$ and remembering that the inequality is reversed when
$r<0$, we obtain
\begin{equation}\label{eq:chain-pos}
\Phi(\chi r,h',w_\varsigma)\ \le\ \Big(\frac{C(1+|r|)}{h''-h'}\Big)^{2/r}\Phi(r,h'',w_\varsigma)
\qquad (0<r<1),
\end{equation}
\begin{equation}\label{eq:chain-neg}
\Phi(\chi r,h',w_\varsigma)\ \ge\ \Big(\frac{C(1+|r|)}{h''-h'}\Big)^{2/r}\Phi(r,h'',w_\varsigma)
\qquad (r<0).
\end{equation}

\emph{Step 3: iteration towards the infimum.} Fix $r_0>0$ (to be chosen in Step 4) and set
\[
r_k:=-r_0\chi^{k},\qquad h_k:=\delta\Big(1+\tfrac32\,2^{-k}\Big),\qquad k\ge0 ,
\]
so that $r_k\to-\infty$, $h_0=\frac{5\delta}{2}$, $h_k\downarrow\delta$ and
$h_k-h_{k+1}=\frac{3\delta}{2}2^{-(k+1)}$. Applying \eqref{eq:chain-neg} with $r=r_k$,
$h'=h_{k+1}$, $h''=h_k$ gives
\[
\Phi(r_{k+1},h_{k+1},w_\varsigma)\ \ge\
\Big(\frac{2^{k+2}C\,(1+r_0\chi^{k})}{3\delta}\Big)^{-\frac{2}{r_0\chi^{k}}}\Phi(r_k,h_k,w_\varsigma) ,
\]
where we used $|r_k|=r_0\chi^{k}$ and $h_k-h_{k+1}=\frac{3\delta}{2}2^{-(k+1)}$. Since
$1+r_0\chi^{k}\le(1+r_0^{-1})r_0\chi^{k}$, taking logarithms and summing over $k\ge0$ leads to the
series
\[
\sum_{k\ge0}\frac{2}{r_0\chi^{k}}\Big[\log\tfrac{4C(1+r_0^{-1})r_0}{3\delta}+k\log 2+k\log\chi\Big] ,
\]
which converges because $\chi>1$. Hence there is $C_1=C_1(r_0,\chi,\delta,C)>0$, independent of
$k$, such that $\Phi(r_k,h_k,w_\varsigma)\ge C_1\,\Phi(-r_0,\frac{5\delta}{2},w_\varsigma)$ for every
$k$. On the other hand $h_k\ge\delta$ and $r_k<0$, so that
$\Phi(r_k,h_k,w_\varsigma)\le\Phi(r_k,\delta,w_\varsigma)$, and the latter converges to
$\Phi(-\infty,\delta,w_\varsigma)$ as $k\to\infty$. Therefore
\begin{equation}\label{eq:toinf}
\inf_{B_{\delta}(x)}w_\varsigma=\Phi(-\infty,\delta,w_\varsigma)\ \ge\ C_1\,\Phi(-r_0,\tfrac{5\delta}{2},w_\varsigma).
\end{equation}

\emph{Step 4: Trudinger's exponential estimate.} We claim that there exist $r_0>0$ and $C_2>0$,
both independent of $\varsigma$, such that
\begin{equation}\label{eq:trudinger}
\Phi\big(r_0,\tfrac{5\delta}{2},w_\varsigma\big)\ \le\ C_2\,\Phi\big(-r_0,\tfrac{5\delta}{2},w_\varsigma\big).
\end{equation}
Put
\[
k_\varsigma:=\exp\Big(\frac{1}{|B_{5\delta}(x)|}\int_{B_{5\delta}(x)}\log w_\varsigma\,dy\Big),
\qquad \widetilde w:=\log\frac{w_\varsigma}{k_\varsigma},
\]
so that $\widetilde w$ has zero \emph{Lebesgue} mean on $B_{5\delta}(x)$. We stress that
$w_\varsigma$ itself is \emph{not} renormalized, it is the function $v-u+\varsigma$, which is the
one satisfying \eqref{eq:scp-hyp}, and that only the auxiliary function $\widetilde w$ carries the
constant $k_\varsigma$; this is legitimate because
$D\widetilde w=Dw_\varsigma/w_\varsigma$ irrespective of $k_\varsigma$, and because $k_\varsigma$
will cancel identically in the final display.

Choosing in \eqref{eq:cacc-log} a cut-off $\eta$ with $\eta\equiv1$ on $B_{5\delta}(x)$ and
$\supp\eta\subset B_{6\delta}(x)$, this is the only point where the hypothesis
$\overline{B_{6\delta}(x)}\subset\Omega'$ is needed, we get from \eqref{eq:fundamental-log}
\begin{equation}\label{eq:logenergy}
\int_{B_{5\delta}(x)}\rho\,|D\widetilde w|^2dy\ \le\ C_3 ,
\end{equation}
with $C_3$ independent of $\varsigma$. Since $\overline{\widetilde w}_{B_{5\delta}(x)}=0$, the
Poincar\'e--Sobolev inequality \eqref{eq:sobolev-mean} and \eqref{eq:logenergy} give
\begin{equation}\label{eq:logLq}
\Phi\big(q,5\delta,\widetilde w\big)=\|\widetilde w\|_{L^{q}(B_{5\delta}(x),\mu)}\ \le\ C_4 ,
\end{equation}
again with $C_4$ independent of $\varsigma$. This uniformity is what allows us to let
$\varsigma\to0$ at the end.

Next we derive a Caccioppoli inequality for powers of $\widetilde w$. Let $\beta\ge1$ and let
$\eta\in C_c^1(B_{5\delta}(x))$, $0\le\eta\le1$. Using in \eqref{eq:scp-hyp} the (nonnegative) test
function
\[
\phi:=\eta^2\,\frac{\Psi(\widetilde w)}{w_\varsigma},\qquad \Psi(t):=|t|^{\beta}+(2\beta)^{\beta},
\]
whose gradient is
\[
D\phi=\frac{2\eta\,\Psi(\widetilde w)}{w_\varsigma}D\eta
+\frac{\eta^2}{w_\varsigma^{2}}\Big[\beta\operatorname{sgn}(\widetilde w)|\widetilde w|^{\beta-1}-\Psi(\widetilde w)\Big]Dw_\varsigma ,
\]
and recalling $D\widetilde w=Dw_\varsigma/w_\varsigma$, $(F,Dw_\varsigma)\ge\hat C\rho|Dw_\varsigma|^2$
and $|F|\le\check C\rho|Dw_\varsigma|$, we obtain
\begin{align*}
&\int\frac{\eta^2}{w_\varsigma^{2}}\Big[\Psi(\widetilde w)-\beta\operatorname{sgn}(\widetilde w)|\widetilde w|^{\beta-1}\Big](F,Dw_\varsigma)\,dy\\
&\qquad\le\ 2\check C\int\eta\,\Psi(\widetilde w)\,\rho\,|D\widetilde w||D\eta|\,dy
+|\Lambda|\int\eta^2\Psi(\widetilde w)\,dy ,
\end{align*}
where we used $0\le v-u\le w_\varsigma$ in the last term. For $\beta\ge1$, Young's inequality with
exponents $\frac{\beta}{\beta-1}$ and $\beta$ yields
$2\beta|t|^{\beta-1}\le\frac{\beta-1}{\beta}|t|^{\beta}+\frac1\beta(2\beta)^{\beta}\le\Psi(t)$,
whence
\begin{equation}\label{eq:elem}
\Psi(t)-\beta\operatorname{sgn}(t)|t|^{\beta-1}\ \ge\ \Psi(t)-\beta|t|^{\beta-1}\ \ge\ \beta|t|^{\beta-1}\ \ge\ 0 .
\end{equation}
Consequently
\[
\beta\hat C\int\eta^2\rho\,|\widetilde w|^{\beta-1}|D\widetilde w|^2
\le 2\check C\int\eta\rho\,\big(|\widetilde w|^{\beta}+(2\beta)^{\beta}\big)|D\widetilde w||D\eta|
+|\Lambda|\int\eta^2\Psi(\widetilde w) .
\]
We split the first integral on the right. Young's inequality gives
\[
2\check C\int\eta\rho|\widetilde w|^{\beta}|D\widetilde w||D\eta|
\le\frac{\beta\hat C}{2}\int\eta^2\rho|\widetilde w|^{\beta-1}|D\widetilde w|^2
+C\int\rho\,|\widetilde w|^{\beta+1}|D\eta|^2 ,
\]
while, for the second piece, \eqref{eq:logenergy} gives
\[
2\check C(2\beta)^{\beta}\int\eta\rho|D\widetilde w||D\eta|
\le (2\beta)^{\beta}\Big[\check C\int\eta^2\rho|D\widetilde w|^2+\check C\int\rho|D\eta|^2\Big]
\le C(2\beta)^{\beta}+\check C(2\beta)^\beta\int\rho|D\eta|^2 .
\]
Finally $\int\eta^2\Psi(\widetilde w)\le C\int\eta^2\big(|\widetilde w|^{\beta+1}+(2\beta)^{\beta}\big)$
because $|t|^{\beta}\le|t|^{\beta+1}+1$ and $1\le(2\beta)^{\beta}$, and, since
$\int\eta^2 dy\ge|B_{5\delta/2}(x)|$ for all the cut-offs used below, the additive constant
$C(2\beta)^{\beta}$ can be written as $C(\delta)\int\eta^2(2\beta)^{\beta}dy$. Absorbing and
collecting terms we arrive at
\begin{equation}\label{eq:cacc-log-beta}
\int\rho\,\eta^2|\widetilde w|^{\beta-1}|D\widetilde w|^2dy\ \le\
C\int\big(|\widetilde w|^{\beta+1}+(2\beta)^{\beta}\big)\big(\eta^2+\rho|D\eta|^2\big)dy ,
\end{equation}
which is the analogue of \eqref{eq:cacc-comp} for the logarithm, with the extra additive term
$(2\beta)^{\beta}$.

Now set $r:=\beta+1\ge2$ and $\widehat w:=|\widetilde w|^{r/2}$, so that
$|D\widehat w|^2=\frac{r^2}{4}|\widetilde w|^{\beta-1}|D\widetilde w|^2$. Exactly as in Step 1,
combining \eqref{eq:cacc-log-beta} with Theorem \ref{thm:sobolev} and \eqref{eq:passtomu} we get,
for $\frac{5\delta}{2}\le h'<h''\le5\delta$,
\[
\Phi(\chi r,h',\widetilde w)^{r}\ \le\ \Big(\frac{C(1+|r|)}{h''-h'}\Big)^{2}
\Big[\Phi(r,h'',\widetilde w)^{r}+(2\beta)^{\beta}\mu(B_{5\delta}(x))\Big] .
\]
Note that here $r\ge2$, so that $1+|r|\le\frac32 r$.
Taking the power $1/r\le\frac12$ and using $(a+b)^{1/r}\le a^{1/r}+b^{1/r}$ together with
$\big((2\beta)^{\beta}\mu(B_{5\delta}(x))\big)^{1/r}\le(2\beta)^{\frac{\beta}{\beta+1}}\max\{1,\mu(B_{5\delta}(x))\}
\le 2r\max\{1,\mu(B_{5\delta}(x))\}$, we obtain
\begin{equation}\label{eq:chain-log}
\Phi(\chi r,h',\widetilde w)\ \le\ \Big(\frac{C|r|}{h''-h'}\Big)^{2/r}
\Big[\Phi(r,h'',\widetilde w)+\gamma\,r\Big],\qquad r\ge2 ,
\end{equation}
with $\gamma=\gamma(\mu(B_{5\delta}(x)))>0$.

We iterate \eqref{eq:chain-log} along $\mu_k:=\frac{5\delta}{2}(1+2^{-k})$ (so that
$\mu_0=5\delta$, $\mu_k\downarrow\frac{5\delta}{2}$ and $\mu_{k-1}-\mu_k=\frac{5\delta}{2}2^{-k}$)
and $r=\chi^{k-1}q$. Setting $A_k:=\Phi(\chi^{k}q,\mu_k,\widetilde w)$ and
$\theta_k:=\big(\tfrac{2^{k+1}C\chi^{k-1}q}{5\delta}\big)^{\frac{2}{\chi^{k-1}q}}$, inequality
\eqref{eq:chain-log} reads $A_k\le\theta_k\,(A_{k-1}+\gamma\chi^{k-1}q)$. Since
$\sum_{k\ge1}\frac{2}{\chi^{k-1}q}\log\big(\tfrac{2^{k+1}C\chi^{k-1}q}{5\delta}\big)<\infty$, the
infinite product $\Theta:=\prod_{k\ge1}\theta_k$ converges, and an elementary induction gives
\[
A_k\ \le\ \Theta A_0+\gamma\Theta q\sum_{j=1}^{k}\chi^{\,j-1}\ \le\ C\big(A_0+\chi^{k}q\big) ,
\]
with $C$ independent of $k$. Given $m\ge q$, let $k_m:=\min\{h\in\N:\ \chi^{h}q\ge m\}$, so that
$\chi^{k_m}q\le\chi m$; by \eqref{eq:holderPhi},
\begin{equation}\label{eq:momentbound}
\Phi\big(m,\tfrac{5\delta}{2},\widetilde w\big)\ \le\
C\,\Phi\big(\chi^{k_m}q,h_{k_m},\widetilde w\big)\ \le\ C\big(\Phi(q,5\delta,\widetilde w)+m\big)
\ \le\ C_5\,(1+m),
\end{equation}
where in the last step we used \eqref{eq:logLq}; the constant $C_5$ does not depend on $\varsigma$
nor on $m$.

We can now conclude Step 4. Note first that \eqref{eq:momentbound} also holds, after enlarging
$C_5$, for the finitely many integers $1\le k<q$: indeed \eqref{eq:holderPhi} gives
$\Phi(k,\frac{5\delta}{2},\widetilde w)\le\mu(B_{5\delta/2}(x))^{\frac1k-\frac1q}\Phi(q,\frac{5\delta}{2},\widetilde w)$,
which is bounded by \eqref{eq:logLq}. Expanding the exponential and using \eqref{eq:momentbound}
for every $k\ge1$ (the term $k=0$ equals $\mu(B_{5\delta/2}(x))$ and is absorbed into the
constant $C$ below),
\[
\int_{B_{5\delta/2}(x)}e^{r_0|\widetilde w|}d\mu
\le\sum_{k\ge0}\frac{r_0^{k}}{k!}\int_{B_{5\delta/2}(x)}|\widetilde w|^{k}d\mu
=\sum_{k\ge0}\frac{r_0^{k}\,\Phi\big(k,\frac{5\delta}{2},\widetilde w\big)^{k}}{k!}
\le C\sum_{k\ge0}\frac{\big(C_5 r_0\big)^{k}(1+k)^{k}}{k!} .
\]
Since $\frac{(1+k)^{k}}{k!}\le e^{k+1}$ by Stirling's bound $k!\ge k^{k}e^{-k}$, the last series
converges as soon as $r_0<\frac{1}{eC_5}$; we fix such an $r_0$, which depends only on the data.
Hence $\int_{B_{5\delta/2}(x)}e^{r_0|\widetilde w|}d\mu\le C$ and therefore
\[
\int_{B_{5\delta/2}(x)}w_\varsigma^{\,r_0}d\mu\cdot\int_{B_{5\delta/2}(x)}w_\varsigma^{-r_0}d\mu
=k_\varsigma^{\,r_0}\!\int e^{r_0\widetilde w}d\mu\cdot k_\varsigma^{-r_0}\!\int e^{-r_0\widetilde w}d\mu
\le\Big(\int e^{r_0|\widetilde w|}d\mu\Big)^{2}\le C^{2} ,
\]
the constant $k_\varsigma$ cancelling exactly, as announced. Raising to the power $\frac1{r_0}$
gives \eqref{eq:trudinger}, with a constant $C_2=C^{2/r_0}$ independent of $\varsigma$.

\emph{Step 5: conclusion.} Let $0<s<\chi$. If $s\le r_0$, then by \eqref{eq:holderPhi}
\[
\Phi(s,2\delta,w_\varsigma)\le\Phi\big(s,\tfrac{5\delta}{2},w_\varsigma\big)
\le\mu\big(B_{5\delta/2}(x)\big)^{\frac1s-\frac1{r_0}}\Phi\big(r_0,\tfrac{5\delta}{2},w_\varsigma\big).
\]
If instead $r_0<s<\chi$, choose $k_0\in\N$ so large that $r_1:=s\chi^{-(k_0+1)}\le r_0$ and set
$\hat r_k:=r_1\chi^{k}$ and $\frac{5\delta}{2}=\varrho_0>\varrho_1>\dots>\varrho_{k_0+1}=2\delta$. Since
$\hat r_{k}\le r_1\chi^{k_0}=s/\chi<1$ for $k\le k_0$, inequality \eqref{eq:chain-pos} may be applied at
each of the $k_0+1$ steps, and it yields
$\Phi(s,2\delta,w_\varsigma)\le C\,\Phi\big(r_1,\frac{5\delta}{2},w_\varsigma\big)$; here $C$ is the
product of the $k_0+1$ factors appearing in \eqref{eq:chain-pos}, hence a finite constant depending
on $s$, on $\delta$ and on the data, but not on $\varsigma$ (this is the only place where the
constant of Theorem \ref{thm:abstract} depends on $s$). By
\eqref{eq:holderPhi} and \eqref{eq:trudinger},
\[
\Phi\big(r_1,\tfrac{5\delta}{2},w_\varsigma\big)\le C\,\Phi\big(r_0,\tfrac{5\delta}{2},w_\varsigma\big)
\le C\,\Phi\big(-r_0,\tfrac{5\delta}{2},w_\varsigma\big) .
\]
In both cases, combining with \eqref{eq:toinf} we obtain
\[
\Phi(s,2\delta,w_\varsigma)\ \le\ C\,\Phi\big(-r_0,\tfrac{5\delta}{2},w_\varsigma\big)
\ \le\ C\,\inf_{B_{\delta}(x)}w_\varsigma ,
\]
with $C$ independent of $\varsigma$. Letting $\varsigma\downarrow0$ and using monotone convergence
on the left and $\inf_{B_{\delta}(x)}w_\varsigma=\varsigma+\inf_{B_{\delta}(x)}w$ on the right, we
conclude that
\[
\Big(\int_{B_{2\delta}(x)}(v-u)^{s}\,d\mu\Big)^{1/s}\ \le\ C\,\inf_{B_{\delta}(x)}(v-u)
\qquad\text{for every }0<s<\chi .
\]

\emph{Step 6: the strong comparison principle.} Let $\Omega_0\subseteq\Omega'$ be a connected
subdomain and set $U:=\{y\in\Omega_0:\ u(y)=v(y)\}$. Since $u,v$ are continuous, $U$ is closed in
$\Omega_0$. If $y_0\in U$, choose $\delta>0$ with $\overline{B_{6\delta}(y_0)}\subset\Omega_0$; then
$\inf_{B_{\delta}(y_0)}(v-u)=0$, hence $\int_{B_{2\delta}(y_0)}(v-u)^s\,d\mu=0$ and, since $\mu$ and $dx$
are mutually absolutely continuous, $u\equiv v$ on $B_{2\delta}(y_0)$. Thus $U$ is open, and the
connectedness of $\Omega_0$ gives the alternative.
\end{proof}

\end{document}
